\begin{table}[!h]
\caption{Ti 在不同状态下的表现形式。}

\resizebox{\linewidth}{!}{%

\begin{tabular}{ccccccc}
\toprule 
\textbf{功能} & \textbf{角色} & \textbf{状态} & \textbf{涉及类型} & \textbf{触发器} & \textbf{关键字/核心台词} & \textbf{表现形式 ｜ 例子}\tabularnewline
\midrule
\midrule 
\multirow{12}{*}{Ti} & Hero/Parent & 正常 & - & N/A & 逻辑，解构，定义 & 追求真理，刨根问底，无视权威\tabularnewline
\cmidrule{2-7}
 & \multirow{2}{*}{Child} & 游戏(Playfulness) & INFJ, ISFJ & N/A & \multirow{2}{*}{居然是这样！} & 纯粹好奇，钻研冷门，解谜乐趣\tabularnewline
\cmidrule{3-5}\cmidrule{7-7}
 &  & 幼稚(Immaturity) & INFJ, ISFJ & 压力 &  & 逻辑死板，钻牛角尖，固执己见\tabularnewline
\cmidrule{2-7}
 & \multirow{3}{*}{Anima} & 正常(Aspiration) & ENFJ, ESFJ & N/A & “我想拥有智慧” & \begin{cellvarwidth}[t]
\centering
引经据典，崇拜专家

读艰深晦涩的书，引用名人名言
\end{cellvarwidth}\tabularnewline
\cmidrule{3-7}
 &  & 压抑(Suppression) & ENFJ, ESFJ & 防御性屏蔽 & “质疑是 Fe 的破坏” & \begin{cellvarwidth}[t]
\centering
智力依赖，盲从共识

“大家都这么说”，拒绝独立思考，屏蔽异见
\end{cellvarwidth}\tabularnewline
\cmidrule{3-7}
 &  & 劫持/抓取(Grip) & ENFJ, ESFJ & 人际背叛，缺乏反馈 & “人类是愚蠢且无可救药的” & 冷酷的犬儒主义，用刺耳的逻辑攻击他人软肋\tabularnewline
\cmidrule{2-7}
 & \multirow{2}{*}{Nemesis} & 防御表现(Stubbornness) & ENTJ, ESTJ & 主导功能受到挑战 & \multirow{2}{*}{“这没意义”} & \begin{cellvarwidth}[t]
\centering
傲慢驳斥：若是不能马上应用，

就拒绝深入思考逻辑细节，认为那是浪费时间
\end{cellvarwidth}\tabularnewline
\cmidrule{3-5}\cmidrule{7-7}
 &  & 投射怀疑(Paranoia) & ENTJ, ESTJ & 对他人该功能莫名厌烦 &  & \begin{cellvarwidth}[t]
\centering
怀疑智商：认为别人的精细逻辑是“书呆子气”，

试图用权威压倒道理
\end{cellvarwidth}\tabularnewline
\cmidrule{2-7}
 & \multirow{2}{*}{Critic} & 对外打压 & INTJ, ISTJ & 防御机制 & \multirow{2}{*}{批判逻辑: 你真蠢} & 智力羞辱，冷漠指出他人逻辑漏洞\tabularnewline
\cmidrule{3-5}\cmidrule{7-7}
 &  & 对内自毁 & INTJ, ISTJ & 防御机制 &  & 自我怀疑，质疑自己不够聪明\tabularnewline
\cmidrule{2-7}
 & Trickster & 忽视/贬低 & ENFP, ESFP & 无意识的黑洞 & “太较真了没意思，差不多就行了。” & \begin{cellvarwidth}[t]
\centering
逻辑疏松，讨厌定义和辩论

经常自我矛盾却不自知，只关注“好不好”
\end{cellvarwidth}\tabularnewline
\cmidrule{2-7}
 & Demon & 毁灭一切 & INFP, ISFP & 彻底崩溃，严重侵犯，极度绝望 & 诛心真话 | Cold Truth & 最恶毒的逻辑穿刺，摧毁对方自尊\tabularnewline
\bottomrule
\end{tabular}

}
\end{table}
